\subsubsection{Window Functions}

The signal for a window function is a per-row result that still depends on
other rows, such as a rank, a neighbor value, or a running total. Where
\texttt{GROUP BY} would collapse the group to one row, a window function keeps
every row and appends the computed column.

Every window function shares the same shape:
\texttt{f() OVER (PARTITION BY ... ORDER BY ...)}. \texttt{PARTITION BY} splits
the rows into independent groups, \texttt{ORDER BY} sets the order within each
group, and the function is evaluated per row. All of them are computed at the
\texttt{SELECT} step, so to filter on their output you wrap them in a subquery
or CTE.

\paragraph{Rank-then-filter with ROW\_NUMBER().} Use this for top-$N$ per
group when you need to keep the full row, since \texttt{GROUP BY} would collapse
the other columns away. Number rows inside each partition, then filter on that
number in an outer query.

Choose the ranker by tie policy: \texttt{ROW\_NUMBER} assigns distinct numbers
and breaks ties arbitrarily, \texttt{RANK} repeats a rank and leaves gaps,
\texttt{DENSE\_RANK} repeats with no gaps.

\begin{lstlisting}[style=sql, caption={Highest-paid employee per department}]
SELECT department, name, salary
FROM (
    SELECT department, name, salary,
           ROW_NUMBER() OVER (PARTITION BY department
                              ORDER BY salary DESC) AS rn
    FROM employees
) ranked
WHERE rn = 1;
\end{lstlisting}

\vspace{-4ex}

\paragraph{Offset access with LAG() and LEAD().} The signal is a comparison
against the previous or next row: a delta, a period-over-period change, or a
return. These read a value from another row relative to the current one,
without a self-join. \texttt{LAG(col, k, default)} returns the value \texttt{k}
rows earlier in the ordered partition, \texttt{LEAD} looks \texttt{k} rows
ahead. Both default to \texttt{k = 1} and return \texttt{NULL} past the
partition edge unless a fill \texttt{default} is supplied. The \texttt{ORDER BY}
inside \texttt{OVER} defines what ``previous'' means and is mandatory.
\texttt{PARTITION BY} resets the sequence at each group boundary, so a lag never
bleeds from one partition into the next.

\begin{lstlisting}[style=sql, caption={Daily returns per ticker with LAG}]
SELECT ticker, trade_date, close,
       close / LAG(close) OVER (PARTITION BY ticker
                                ORDER BY trade_date) - 1 AS daily_return
FROM prices;
\end{lstlisting}

\vspace{-1ex}

The first row of each partition has no predecessor, so its \texttt{LAG} and the
derived return are \texttt{NULL}, which is correct since there is no prior price
to compare against. \texttt{LEAD} is the same tool aimed forward, used below for
gap detection.

\paragraph{Running and rolling aggregates.} The signal is ``running'',
``cumulative'', or ``moving average''. An ordinary aggregate gains a window when
you attach \texttt{OVER (... ORDER BY ...)}. With an \texttt{ORDER BY} and no
explicit frame, the frame defaults to everything from the partition start up to
the current row, giving a running total. Add an explicit \texttt{ROWS BETWEEN}
frame to get a fixed moving window instead.

\begin{lstlisting}[style=sql, caption={20-day moving average and running volume}]
SELECT ticker, trade_date, close,
       AVG(close) OVER (PARTITION BY ticker ORDER BY trade_date
                        ROWS BETWEEN 19 PRECEDING AND CURRENT ROW) AS ma_20,
       SUM(volume) OVER (PARTITION BY ticker ORDER BY trade_date)  AS cum_volume
FROM prices;
\end{lstlisting}

\vspace{-1ex}

\texttt{ma\_20} is the current row plus the 19 before it, and \texttt{cum\_volume}
uses the default frame for a running total. These are the SQL forms of a pandas
\texttt{.rolling(20).mean()} and \texttt{.expanding().sum()}.

One subtlety: the default frame is \texttt{RANGE UNBOUNDED PRECEDING}, which
includes every row tied on the \texttt{ORDER BY} key, not just rows physically
up to the current one. When ties on the order key exist and you want a
row-precise window, state \texttt{ROWS BETWEEN} explicitly.

\paragraph{Gaps and islands.} The signal is any mention of consecutive values,
streaks, or runs. An \emph{island} is a maximal run of consecutive values. The
trick: subtract a dense row number from the value. Across a consecutive run
both sides increase by one per step, so the difference stays constant, and it
jumps exactly at a gap. Group by the difference to collapse each run.

\begin{lstlisting}[style=sql, caption={Islands of consecutive ids}]
WITH numbered AS (
    -- DISTINCT first so duplicate ids do not desync the row number
    SELECT DISTINCT log_id,
           ROW_NUMBER() OVER (ORDER BY log_id) AS rn
    FROM logs
)
SELECT MIN(log_id) AS island_start,
       MAX(log_id) AS island_end,
       COUNT(*)    AS length
FROM numbered
GROUP BY log_id - rn
ORDER BY island_start;
\end{lstlisting}

The \texttt{DISTINCT} is load-bearing. If \texttt{log\_id} repeats, \texttt{rn}
still advances on the duplicate while the value does not, so
\texttt{log\_id - rn} shifts and splits an island that should stay whole. For
consecutive \emph{dates} rather than integers, offset by an interval:
\texttt{log\_date - rn * INTERVAL `1 day'} is constant within a run.

\begin{lstlisting}[style=sql, caption={Gaps via LEAD}]
SELECT log_id + 1  AS missing_from,
       next_id - 1 AS missing_to
FROM (
    SELECT log_id,
           LEAD(log_id) OVER (ORDER BY log_id) AS next_id
    FROM logs
) t
WHERE next_id - log_id > 1;   -- values missing between two present ids
\end{lstlisting}
